\documentclass[11pt,a4paper]{article}

% Encodage et polices
\usepackage[utf8]{inputenc}
\usepackage[T1]{fontenc}
\usepackage[french]{babel}

% Mise en page
\usepackage{geometry}
\geometry{a4paper, margin=2.5cm}

% Mathématiques avancées
\usepackage{amsmath,amssymb,amsfonts}

% Graphiques et tableaux
\usepackage{graphicx}
\usepackage{booktabs}
\usepackage{array}

% Liens hypertextes
\usepackage{hyperref}
\hypersetup{
    colorlinks=true,
    linkcolor=blue,
    filecolor=magenta,      
    urlcolor=cyan,
    citecolor=blue
}

% Texte barré et souligné
\usepackage[normalem]{ulem}

% Blocs de code
\usepackage{listings}
\usepackage{xcolor}

\definecolor{codebg}{rgb}{0.96,0.96,0.96}
\definecolor{codecomment}{rgb}{0.25,0.5,0.35}
\definecolor{codekeyword}{rgb}{0.13,0.13,1}
\definecolor{codestring}{rgb}{0.8,0.1,0.1}

\lstdefinestyle{moderncode}{
    backgroundcolor=\color{codebg},
    commentstyle=\color{codecomment}\textit,
    keywordstyle=\color{codekeyword}\textbf,
    numberstyle=\tiny\color{gray},
    stringstyle=\color{codestring},
    basicstyle=\ttfamily\footnotesize,
    breakatwhitespace=false,
    breaklines=true,
    captionpos=b,
    keepspaces=true,
    numbers=left,
    numbersep=8pt,
    showspaces=false,
    showstringspaces=false,
    showtabs=false,
    tabsize=4,
    frame=single,
    rulecolor=\color{black!20},
    literate=
      {á}{{\'a}}1 {é}{{\'e}}1 {í}{{\'i}}1 {ó}{{\'o}}1 {ú}{{\'u}}1
      {Á}{{\'A}}1 {É}{{\'E}}1 {Í}{{\'I}}1 {Ó}{{\'O}}1 {Ú}{{\'U}}1
      {à}{{\`a}}1 {è}{{\`e}}1 {ì}{{\`i}}1 {ò}{{\`o}}1 {ù}{{\`u}}1
      {À}{{\`A}}1 {È}{{\'E}}1 {Ì}{{\`I}}1 {Ò}{{\`O}}1 {Ù}{{\`U}}1
      {ä}{{\"a}}1 {ë}{{\"e}}1 {ï}{{\"i}}1 {ö}{{\"o}}1 {ü}{{\"u}}1
      {Ä}{{\"A}}1 {Ë}{{\"E}}1 {Ï}{{\"I}}1 {Ö}{{\"O}}1 {Ü}{{\"U}}1
      {â}{{\^a}}1 {ê}{{\^e}}1 {î}{{\^i}}1 {ô}{{\^o}}1 {û}{{\^u}}1
      {Â}{{\^A}}1 {Ê}{{\^E}}1 {Î}{{\^I}}1 {Ô}{{\^O}}1 {Û}{{\^U}}1
      {œ}{{\oe}}1 {Œ}{{\OE}}1 {æ}{{\ae}}1 {Æ}{{\AE}}1 {ç}{{\c c}}1
      {Ç}{{\c C}}1 {«}{{\guillemotleft}}1 {»}{{\guillemotright}}1
}
\lstset{style=moderncode}

\title{Rapport de Démonstration : Conversion Markdown vers LaTeX}
\author{}
\date{\today}

\begin{document}

\maketitle
\vspace{1em}


Bienvenue dans ce document d'exemple pour le projet \textbf{md2tex}.

\subsection{1. Formatage de texte et caractères spéciaux}

On peut mettre du texte en \textbf{gras}, en \textit{italique}, en \textbf{\textit{gras et italique}}, ou encore en \sout{texte barré}.
Les caractères spéciaux comme 100\%, les esperluettes \& cie, ainsi que les variables avec underscores comme \texttt{taux\_variation\_annuel} sont convenablement échappés pour éviter les erreurs de compilation LaTeX.

\subsection{2. Formules Mathématiques}

\subsubsection{Formule en ligne}
Le théorème de Pythagore s'énonce $a^2 + b^2 = c^2$.
La célèbre équation d'Einstein est $E = m \cdot c^2$.
Pour une variable aléatoire $X \sim \mathcal{N}(\mu, \sigma^2)$, la densité est donnée par $f(x)$.

\subsubsection{Formule en bloc}
Voici l'intégrale de Gauss :

\begin{equation*}
\int_{-\infty}^{+\infty} e^{-x^2} dx = \sqrt{\pi}
\end{equation*}

Et la formule de décomposition en série de Fourier :

\begin{equation*}
f(x) = \frac{a_0}{2} + \sum_{n=1}^{\infty} \left( a_n \cos(nx) + b_n \sin(nx) \right)
\end{equation*}

\noindent\rule{\linewidth}{0.4pt}

\subsection{3. Listes}

\subsubsection{Liste non ordonnée}
\begin{itemize}
  \item Premier élément important
  \item Deuxième élément avec une sous-liste :
\begin{itemize}
  \item Sous-élément A
  \item Sous-élément B avec formule : $x_k \in \mathbb{R}$
\end{itemize}
  \item Troisième élément
\end{itemize}

\subsubsection{Liste ordonnée}
\begin{enumerate}
  \item Étape 1 : Initialisation des variables
  \item Étape 2 : Traitement algorithmique
  \item Étape 3 : Exportation des résultats
\end{enumerate}

\noindent\rule{\linewidth}{0.4pt}

\subsection{4. Blocs de Code}

Voici un extrait de code Python avec coloration syntaxique :

\begin{lstlisting}[language=python]
def fibonacci(n: int) -> list[int]:
    """Génère la suite de Fibonacci jusqu'à n termes."""
    sequence = [0, 1]
    while len(sequence) < n:
        sequence.append(sequence[-1] + sequence[-2])
    return sequence[:n]

if __name__ == "__main__":
    print(fibonacci(10))
\end{lstlisting}

\noindent\rule{\linewidth}{0.4pt}

\subsection{5. Tableaux}

\begin{table}[htbp]
\centering
\begin{tabular}{|l|c|r|c|}
\hline
\textbf{Nom de l'algorithme} & \textbf{Complexité temporelle} & \textbf{Complexité spatiale} & \textbf{Type} \\ \hline
Tri Rapide (QuickSort) & $\mathcal{O}(n \log n)$ & $\mathcal{O}(\log n)$ & Diviser pour régner \\ \hline
Tri par Fusion (MergeSort) & $\mathcal{O}(n \log n)$ & $\mathcal{O}(n)$ & Stable \\ \hline
Tri à Bulles (BubbleSort) & $\mathcal{O}(n^2)$ & $\mathcal{O}(1)$ & Naïf \\ \hline
\end{tabular}
\end{table}

\noindent\rule{\linewidth}{0.4pt}

\subsection{6. Citations et Liens}

\begin{quote}
La science est ce que l'on comprend bien assez pour l'expliquer à un ordinateur. L'art est tout le reste. — Donald Knuth
\end{quote}

Pour en savoir plus, consultez la documentation officielle de \href{https://www.latex-project.org}{LaTeX Project} ou le dépôt du projet.

\end{document}
